\documentclass[11pt,a4paper]{article}
\usepackage[T1]{fontenc}
\usepackage[utf8]{inputenc}
\usepackage[french]{babel}
\usepackage{lmodern}
\usepackage{microtype}
\usepackage[a4paper,margin=2.2cm]{geometry}
\usepackage{xcolor}
\usepackage{enumitem}
\usepackage{hyperref}
\definecolor{esilvblue}{HTML}{153A5B}
\hypersetup{colorlinks=true,urlcolor=esilvblue}
\setlist{nosep,leftmargin=*}
\pagestyle{plain}

\begin{document}

\begin{center}
{\Large\bfseries\color{esilvblue} Projet PI\textsuperscript{2} n\textsuperscript{o} 176}\\[0.3em]
{\large Brief de la réunion de lancement}\\[0.2em]
Mercredi 14 octobre 2026, 16 h--16 h 30
\end{center}

\section*{Résultat attendu de la réunion}
\begin{enumerate}
  \item Valider le sous-jacent et la source de données historiques.
  \item Fixer le périmètre : méthodes, maturités et cible couverture/trading.
  \item Clarifier la part attendue de Julia ou OCaml.
  \item Définir le premier prototype et le contenu du suivi du 23 octobre.
  \item Confirmer l'envoi du rapport et du support le 20 octobre.
\end{enumerate}

\section*{Pitch de 45 secondes}
\begin{quote}
Notre objectif est d'extraire une série temporelle de distributions risque-neutre
à partir de prix d'options, d'en comparer les méthodes de calibration, puis d'étudier
leur stabilité et leur prévisibilité avant de tester une application de couverture
ou de trading. Nous proposons de commencer par SPX à horizon proche de 30 jours,
avec un baseline Breeden--Litzenberger en Julia, contrôlé par des tests de
non-arbitrage et de cohérence probabiliste. Nous comparerons ensuite ce baseline à
la méthode Fourier-cosinus et utiliserons un modèle factoriel pour la dynamique.
\end{quote}

\section*{Cinq questions indispensables}
\begin{enumerate}
  \item Avons-nous accès à OptionMetrics/WRDS, Cboe, Bloomberg ou Refinitiv ?
  \item Validez-vous SPX et une échéance constante proche de 30 jours pour le MVP ?
  \item Faut-il implémenter les trois méthodes de calibration ou en approfondir deux ?
  \item Que signifie précisément « une part substantielle en Julia ou OCaml » ?
  \item La cible finale prioritaire est-elle la couverture, le trading, ou les deux ?
\end{enumerate}

\section*{Proposition pour le suivi du 23 octobre}
Bibliographie structurée ; source ou schéma définitif des données ; prototype
Breeden--Litzenberger ; graphiques de smile, calls et densité ; contrôles de
positivité, masse, moyenne/forward et convexité ; protocole de comparaison avec
Fourier ; planning et responsabilités.

\section*{À ne pas promettre}
Une stratégie rentable avant test hors échantillon et coûts ; une RND fiable à
partir de prix bruts ; le TTF sans accès confirmé aux données ICE ; une probabilité
« réelle », puisque la RND incorpore aussi les prix du risque.

\end{document}
